%% toc: P1 8 8.3,8.4,8.7
\chapter{Substitution, Equality Replacement, and Equivalence Signs}
\chaptermark{Substitution and Equivalence Signs}
\label{sec:p01-w02-p01-025-overview}
\index{substitution!and replacement}\index{equality!congruence}
\index{equivalence sign!typed roles}

This chapter collects three pieces of notation that recur throughout the
logical parts of the handbook and that are easy to confuse: the substitution
and replacement operations, the laws that the identity sign obeys in the
natural-deduction calculus, and the several signs written with a triple bar.

\emph{Purpose.}  A reader who meets \(\varphi[t/x]\),
\(\varphi[O\leftarrow\chi]\), \(r=t\), \(\varphi\equiv\psi\),
\(\varphi\equiv_\alpha\psi\), \(\mathcal M\equiv\mathcal N\) or a
definitional sign should be able to find here the defining clause, the type
of the operands, the scope in which that reading is in force, and the
results that connect neighbouring notions.

\emph{Inputs.}  Capture-avoiding substitution, formula complexity and
alpha-equivalence of \cref{sec:fol-syntax-semantics-04} and
\cref{sec:fol-syntax-semantics-03}; first-order satisfaction and logical
equivalence of \cref{sec:fol-semantics-02}; the calculus
\(\mathsf{ND}_{\mathsf{FOL}}\) of \cref{sec:fol-deduction-01}, whose rules
are listed on the registered sheet \texttt{sheet:nd-fol}; the sequent
calculus with equality on \texttt{sheet:lk-eq}; and
the propositional substitution and replacement results of
\cref{sec:p02-propositional-substitution}.

\emph{Conclusions.}  Symmetry, transitivity, term substitutivity, and the
function and relation congruence schemas are proved as separate results of
\(\mathsf{ND}_{\mathsf{FOL}}\)
(\cref{thm:p01-w02-p01-025-symmetry,thm:p01-w02-p01-025-transitivity,thm:p01-w02-p01-025-fun-cong,thm:p01-w02-p01-025-rel-cong,thm:p01-w02-p01-025-term-subst}).
\Cref{thm:p01-w02-p01-025-axiom-metatheorem} shows that the two equality
rules \(=I\) and \(=E\) have exactly the deductive strength of the
reflexivity and congruence axioms over the same calculus without equality
rules.  \Cref{thm:p01-w02-p01-025-lk-symmetric} states the symmetric
replacement rule of the sequent calculus with equality as a derived rule.
\Cref{def:p01-w02-p01-025-equivalence-roles} and
\cref{thm:p01-w02-p01-025-equiv-bridges} fix the typed roles of the
equivalence and definition signs and the valid passages between them.
Every proof is in the annex.

\section{Substitution, replacement, and the propositional exercise}
\label{sec:p01-w02-p01-025-substitution}
\index{substitution!capture-avoiding}\index{replacement!equivalent subformulas}

Two different operations are written with square brackets.  First-order
substitution \(\varphi[t/x]\) replaces the free occurrences of a variable by
a term and renames a binder exactly when a variable of the term would
otherwise be captured; it is defined, with the term clause, the renaming
clause and the deterministic choice of the fresh variable, in
\cref{sec:fol-syntax-semantics-04}.  The same section proves that the
operation preserves complexity (\cref{lem:substitution-complexity}) and
bounds the free variables of the result (\cref{lem:substitution-fv}), shows
by its worked example how the renaming clause keeps an inserted variable
free, and proves that substituting for a variable that is not free yields an
alpha-equivalent formula.  Alpha-equivalence \(\equiv_\alpha\) is the least
constructor congruence generated by renaming one bound variable to a variable
absent from the formula; \cref{cor:alpha-invariance} shows that it preserves
satisfaction.

Propositional uniform substitution \(\varphi[\upsilon]\) and the replacement
\(\varphi[O\leftarrow\chi]\) of selected subformula occurrences are the
operations of \cref{def:p02-uniform-substitution} and
\cref{def:p02-selected-replacement}.  Uniform substitution acts at every
variable occurrence simultaneously; replacement acts at chosen positions of
the parse tree and leaves the other occurrences of the same subformula
unchanged.

\begin{hbconvention}{def:p01-w02-p01-025-exercise-link}{The replacement theorem and its exercise}
\Cref{ex:prop-metatheory-09} asks for a proof that logical equivalence of
\(\mathsf{CPL}\) formulas is an equivalence relation and a congruence for
every primitive connective.  Those are the first two assertions of
\cref{thm:p02-replacement}, whose complete proof is
Annex~\ref{annex:p2.prop.replacement}; the same theorem adds the
selected-occurrence replacement \(\varphi\equiv\varphi[O\leftarrow\chi]\)
for \(\psi\equiv\chi\).  The feedback entry for that exercise in
\cref{app:exercise-feedback} is a solution sketch: it names the decisive
step (each primitive connective is a fixed truth function, so equal input
values give equal output values) and the cases to be checked, namely
negation and each binary connective.  The annex proof carries out exactly
those cases, together with the falsum case and the structural induction
over the formula required for replacement.  A reader who has written a solution can
therefore check it clause by clause against that annex item, and the
uniform-substitution theorem \cref{thm:p02-uniform-substitution}, proved in
Annex~\ref{annex:p2.prop.substitution}, is the companion result for
simultaneous substitution of formulas for variables.
\end{hbconvention}

\section{Equality in natural deduction}
\label{sec:p01-w02-p01-025-equality}
\index{equality!natural deduction}\index{equality!axioms}

\begin{hbconvention}{def:p01-w02-p01-025-nd-conventions}{Calculus conventions for this section}
Fix a first-order signature with function symbols \(\mathcal F\) and
relation symbols \(\mathcal R\), each with an arity \(n\in\mathbb N\); a
nullary function symbol is a constant and a nullary relation symbol is an
atomic formula without arguments.  The identity sign \(=\) is primitive and
logical.  \(\mathsf{ND}_{\mathsf{FOL}}\) is the calculus of
\cref{sec:fol-deduction-01} with exactly the rules of
\texttt{sheet:nd-fol}.  In particular \(=I\) derives \(t=t\) from no
premises, and \(=E\) derives \(\varphi[t/x]\) from the two premises \(r=t\)
and \(\varphi[r/x]\); both instances of \(=E\) use the total capture-avoiding
substitution
and \(\varphi\) may contain any number of free occurrences of \(x\).
Discharge is never vacuous, the eigenvariable of \(\forall I\) is free in no
open assumption, and the eigenvariable \(y\) of \(\exists E\) is absent from
the conclusion, from every other open assumption and from the existential
premise.  As in the source definition of first-order derivability,
\(\Gamma\derives_{\mathsf{ND}_{\mathsf{FOL}}}\varphi\) means that some
derivation of \(\varphi\) has all its open assumptions in \(\Gamma\); the set
\(\Gamma\) may be infinite.
\end{hbconvention}

The bundled statement \cref{thm:p02-equality-laws} lists reflexivity,
symmetry, transitivity and formula replacement together, and the term-model
lemma \cref{lem:equality-congruence} states congruence for the relation
\(\sim_\Delta\) determined by a maximal consistent set.  The following
results state each law separately and without reference to such a set.

\begin{hblemma}{thm:p01-w02-p01-025-symmetry}{proved=annex:p01-w02-p01-025-symmetry}{Symmetry of identity}
For all terms \(r,t\),
\(\{r=t\}\derives_{\mathsf{ND}_{\mathsf{FOL}}}t=r\).
\end{hblemma}

\begin{hblemma}{thm:p01-w02-p01-025-transitivity}{proved=annex:p01-w02-p01-025-transitivity}{Transitivity of identity}
For all terms \(r,t,u\),
\(\{r=t,\ t=u\}\derives_{\mathsf{ND}_{\mathsf{FOL}}}r=u\).
\end{hblemma}

\begin{hblemma}{thm:p01-w02-p01-025-fun-cong}{proved=annex:p01-w02-p01-025-fun-cong}{Function congruence schema}
For every \(n\in\mathbb N\), every \(n\)-ary \(f\in\mathcal F\) and all terms
\(r_1,\ldots,r_n,t_1,\ldots,t_n\),
\[
 \{r_1=t_1,\ldots,r_n=t_n\}
 \derives_{\mathsf{ND}_{\mathsf{FOL}}}
 f(r_1,\ldots,r_n)=f(t_1,\ldots,t_n).
\]
For \(n=0\) the premise set is empty and the conclusion is \(c=c\) for the
constant \(c=f\).  For \(n\geq 1\) the derivation uses one application of
\(=I\) followed by exactly \(n\) applications of \(=E\), one for each
argument place, from left to right.
\end{hblemma}

\begin{hblemma}{thm:p01-w02-p01-025-rel-cong}{proved=annex:p01-w02-p01-025-rel-cong}{Relation congruence schema}
For every \(n\in\mathbb N\), every \(n\)-ary \(R\in\mathcal R\) and all terms
\(r_1,\ldots,r_n,t_1,\ldots,t_n\),
\[
 \{r_1=t_1,\ldots,r_n=t_n,\ R(r_1,\ldots,r_n)\}
 \derives_{\mathsf{ND}_{\mathsf{FOL}}}
 R(t_1,\ldots,t_n).
\]
The same holds with the identity sign in place of a binary \(R\):
\(\{r_1=t_1,\ r_2=t_2,\ r_1=r_2\}\derives_{\mathsf{ND}_{\mathsf{FOL}}}t_1=t_2\).
For \(n=0\) the conclusion is the premise \(R\).  For \(n\geq 1\) the
derivation uses exactly \(n\) applications of \(=E\), one for each argument
place.
\end{hblemma}

\begin{hblemma}{thm:p01-w02-p01-025-term-subst}{proved=annex:p01-w02-p01-025-term-subst}{Substitutivity in terms}
For all terms \(r,t,u\) and every variable \(x\),
\(\{r=t\}\derives_{\mathsf{ND}_{\mathsf{FOL}}}u[r/x]=u[t/x]\).
Substitutivity in formulas is the primitive rule \(=E\) itself.
\end{hblemma}

\Cref{lem:equality-congruence} follows from these lemmas: if
\(\Delta\derives_{\mathsf{ND}_{\mathsf{FOL}}}r_i=t_i\) for each \(i\), then
composing the derivations with
\cref{thm:p01-w02-p01-025-fun-cong,thm:p01-w02-p01-025-rel-cong} gives the
corresponding derivations from \(\Delta\), and a maximal consistent set
contains every formula it derives.

\begin{hbexample}{def:p01-w02-p01-025-cong-example}{A congruence derivation written out}
Let \(f\) be binary and let \(a,b,c,d\) be constants.  The derivation of
\(f(a,b)=f(c,d)\) from \(\{a=c,\ b=d\}\) given by
\cref{thm:p01-w02-p01-025-fun-cong} uses a variable \(w\) and reads
\[
\begin{array}{rll}
1.&f(a,b)=f(a,b)&=I\\
2.&a=c&\text{assumption}\\
3.&f(a,b)=f(c,b)&=E,\ 2,1\ \text{with }\varphi\defeq f(a,b)=f(w,b)\\
4.&b=d&\text{assumption}\\
5.&f(a,b)=f(c,d)&=E,\ 4,3\ \text{with }\varphi\defeq f(a,b)=f(c,w).
\end{array}
\]
At line 3, \(\varphi[a/w]\) is line 1 and \(\varphi[c/w]\) is line 3.  At
line 5, \(\varphi[b/w]\) is line 3 and \(\varphi[d/w]\) is line 5.  The
hole variable \(w\) occurs in none of the four terms, so each instance is
literal replacement of the single occurrence of \(w\).  Had the hole been
written with a variable that already occurs in the terms, the instance
\(\varphi[a/w]\) would replace that occurrence too, and the two displayed
lines would no longer be the premise and conclusion of a single \(=E\) step.
\end{hbexample}

\subsection{Equality axioms against the equality rules}
\label{sec:p01-w02-p01-025-equality-axioms}

\begin{hbdefinition}{def:p01-w02-p01-025-eq-axioms}{Equality axioms and the calculus without equality rules}
For formulas \(e_1,\ldots,e_n,\chi\) write
\(e_1\to(e_2\to\cdots(e_n\to\chi)\cdots)\) for the right-nested
implication; for \(n=0\) it is \(\chi\).  The set
\(\mathrm{Eq}_{\mathcal L}\) of \emph{equality axioms} of the signature
consists of the following sentences, for every choice of pairwise distinct
variables \(x,x_1,\ldots,x_n,y_1,\ldots,y_n\):
\begin{enumerate}
  \item reflexivity: \(\forall x\,(x=x)\);
  \item for each \(n\)-ary \(f\in\mathcal F\), the function congruence axiom
  \[
   \begin{aligned}
   &\forall x_1\cdots\forall x_n\forall y_1\cdots\forall y_n\,
    (x_1=y_1\to(\cdots(x_n=y_n\\
   &\qquad\to f(x_1,\ldots,x_n)=f(y_1,\ldots,y_n))\cdots));
   \end{aligned}
  \]
  \item for each \(n\)-ary \(R\in\mathcal R\), and for \(=\) with \(n=2\),
  the relation congruence axiom
  \[
   \begin{aligned}
   &\forall x_1\cdots\forall x_n\forall y_1\cdots\forall y_n\,
    (x_1=y_1\to(\cdots(x_n=y_n\\
   &\qquad\to(R(x_1,\ldots,x_n)\to R(y_1,\ldots,y_n)))\cdots)).
   \end{aligned}
  \]
\end{enumerate}
For a constant \(c\) the function axiom is \(c=c\); for a nullary \(R\) the
relation axiom is \(R\to R\).  Because every choice of distinct bound
variables is included, \(\mathrm{Eq}_{\mathcal L}\) contains every
alpha-variant of each of its members.

\(\mathsf{ND}^{-}_{\mathsf{FOL}}\) is \(\mathsf{ND}_{\mathsf{FOL}}\) with
the two rules \(=I\) and \(=E\) deleted.  Its language is unchanged, so
\(r=t\) is still a formula, but no rule mentions the identity sign; the
alpha rule and every quantifier rule are retained with their side conditions.
\(\Gamma\derives_{\mathsf{ND}^{-}_{\mathsf{FOL}}}\varphi\) is defined as for
\(\mathsf{ND}_{\mathsf{FOL}}\).
\end{hbdefinition}

Kleene's predicate calculus with equality adds to a Hilbert-style predicate
calculus the axiom \(a=a\) and the schema
\(a=b\supset(A(a)\supset A(b))\), and calls \(a=a\) together with
\(a=b\supset(a=c\supset b=c)\) the equality axioms for \(=\); each
\(n\)-place predicate letter receives \(n\) open equality axioms, one for
each argument place \cite[\S 73, p.~399]{Kleene1952}.  In that calculus the
terms are the variables.  \Cref{def:p01-w02-p01-025-eq-axioms} differs in
three declared respects: it adds a function axiom for every function symbol,
it replaces all argument places in one axiom, and it states the congruence
axiom for \(=\) in the same all-places form as for the other relation
symbols.  The calculus to which the axioms are added is the natural-deduction
calculus of \cref{sec:fol-deduction-01}, not Kleene's Hilbert system.

\begin{hbtheorem}{thm:p01-w02-p01-025-axioms-derivable}{proved=annex:p01-w02-p01-025-axioms-derivable}{The equality rules derive the equality axioms}
Every member of \(\mathrm{Eq}_{\mathcal L}\) is derivable in
\(\mathsf{ND}_{\mathsf{FOL}}\) from no assumptions.
\end{hbtheorem}

\begin{hbtheorem}{thm:p01-w02-p01-025-axiom-replacement}{proved=annex:p01-w02-p01-025-axiom-replacement}{Replacement from the equality axioms}
For every formula \(\varphi\), every variable \(x\) and all terms \(r,t\),
\[
 \mathrm{Eq}_{\mathcal L}\cup\{r=t,\ \varphi[r/x]\}
 \derives_{\mathsf{ND}^{-}_{\mathsf{FOL}}}\varphi[t/x].
\]
Both instances are the capture-avoiding substitution of
\cref{sec:fol-syntax-semantics-04}, including its renaming clause, so the
bound variables of \(\varphi[r/x]\) and \(\varphi[t/x]\) may differ.
\end{hbtheorem}

\begin{hbtheorem}{thm:p01-w02-p01-025-axiom-metatheorem}{proved=annex:p01-w02-p01-025-axiom-metatheorem}{Equality rules and equality axioms have the same strength}
For every set \(\Gamma\) of formulas and every formula \(\varphi\),
\[
 \Gamma\derives_{\mathsf{ND}_{\mathsf{FOL}}}\varphi
 \quad\Longleftrightarrow\quad
 \Gamma\cup\mathrm{Eq}_{\mathcal L}
   \derives_{\mathsf{ND}^{-}_{\mathsf{FOL}}}\varphi.
\]
In particular \(\varphi\) is a theorem of \(\mathsf{ND}_{\mathsf{FOL}}\)
exactly when it is derivable in \(\mathsf{ND}^{-}_{\mathsf{FOL}}\) from the
equality axioms.
\end{hbtheorem}

The left-to-right direction replaces every \(=I\) step by an instance of the
reflexivity axiom and every \(=E\) step by the derivation of
\cref{thm:p01-w02-p01-025-axiom-replacement}; the right-to-left direction
uses \cref{thm:p01-w02-p01-025-axioms-derivable}.  The theorem compares
derivability relations.  It does not say that the two calculi have the same
models: in \(\mathsf{ND}^{-}_{\mathsf{FOL}}\) nothing forces \(=\) to denote
identity, and the axioms only say that it denotes a reflexive relation
preserved by every function and relation symbol of the signature.

\subsection{The sequent calculus with equality}
\label{sec:p01-w02-p01-025-lk}

On \texttt{sheet:lk-eq} equality is governed by the
initial sequent \(\Rightarrow t=t\) and the single left replacement rule \(=L\), which passes from the premise
\(\Gamma,\varphi[r/x]\Rightarrow\Delta\) to the conclusion
\(\Gamma,r=t,\varphi[t/x]\Rightarrow\Delta\).
\Cref{thm:p07-eqlk-replacement} derives the symmetry sequent
\(r=t\Rightarrow t=r\) and both substitution kernels on that sheet, and
\cref{thm:p07-eqlk-congruence} derives transitivity and the function and
relation congruence sequents.  The opposite orientation of \(=L\) is
therefore a derived rule.

\begin{hbcorollary}{thm:p01-w02-p01-025-lk-symmetric}{proved=annex:p01-w02-p01-025-lk-symmetric}{Symmetric replacement in the sequent calculus with equality}
For all finite multisets \(\Gamma,\Delta\), terms \(r,t\), formulas
\(\varphi\) and variables \(x\): if
\(\Gamma,\varphi[t/x]\Rightarrow\Delta\) is derivable on
\texttt{sheet:lk-eq}, then so is
\(\Gamma,r=t,\varphi[r/x]\Rightarrow\Delta\).  The derivation applies \(=L\)
once to the given derivation and then one Cut against a fixed derivation of
\(r=t\Rightarrow t=r\), which itself uses one \(=L\) and one Cut.
\end{hbcorollary}

\begin{hbwarning}{def:p01-w02-p01-025-lk-cut}{The symmetric rule needs Cut}
The derivation of the symmetric rule uses Cut, and Cut cannot be removed
there: \cref{thm:p07-eqlk-cut-failure} shows that \(a=b\Rightarrow b=a\) has
no derivation without Cut and Mix on that sheet, although it is derivable
with them.  The derived rule may therefore be used in derivations that admit
Cut, but it is not available as a step of a cut-free proof search on the
sheet as registered.  In \(\mathsf{ND}_{\mathsf{FOL}}\), by contrast,
symmetry is a two-step derivation from \(=I\) and \(=E\)
(\cref{thm:p01-w02-p01-025-symmetry}); the two calculi derive the same
sequents by \cref{thm:p07-eqlk-nd}, but they do not have the same
cut-free fragments.
\end{hbwarning}

\section{Equivalence and definition signs}
\label{sec:p01-w02-p01-025-signs}
\index{equivalence sign!typed roles}\index{definitional sign}

The editorial defaults of \cref{sec:p01-symbol-registry} are
\(\varphi\Leftrightarrow\psi\) for formula equivalence,
\(\mathcal M\equiv\mathcal N\) for elementary equivalence,
\(a\mathrel{:\equiv}b\) for a definitional sign, and \(\varphi=_\alpha\psi\)
for alpha-equivalence.  Several source presentations keep their own signs for
some of these roles.  The CPL use of \(\varphi\equiv\psi\) is already
recorded in the registry; the convention below names the remaining
source-local uses, each with the type of its operands and its scope.

\begin{hbconvention}{def:p01-w02-p01-025-equivalence-roles}{Typed roles of the triple bar and of the definitional sign}
\begin{enumerate}
  \item \emph{First-order logical equivalence.}  In
  \cref{sec:fol-semantics-02} and in every later first-order source chapter
  that uses it, in particular \cref{sec:fol-cardinality-normalforms-03},
  and in the solutions given in \cref{app:exercise-feedback} for
  \cref{ex:fol-cardinality-normalforms-05} and
  \cref{ex:automated-reasoning-03}, \(\varphi\equiv\psi\) relates two
  formulas of one first-order language \(\mathcal L\): it holds iff every
  \(\mathcal L\)-structure and every assignment satisfy both or neither.  It
  is metalinguistic; it is neither alpha-equivalence nor identity of strings.
  In those two solutions it joins the lines of a displayed chain of
  first-order formulas, each line logically equivalent to the one before.
  \item \emph{Metalinguistic equivalence in displayed propositional
  calculations.}  In \cref{sec:writing-foundations-05}, \(\equiv\) joins the
  lines of a displayed chain of propositional formulas and asserts that
  consecutive lines are logically equivalent.  It is a sign of the
  metalanguage; the object-language biconditional is \(\leftrightarrow\).
  \item \emph{Source definitional sign.}  Inside the restoration scopes that
  reproduce the frozen source, the macro \texttt{\textbackslash defeq} prints
  \(\mathrel{:=}\), as in the source definition
  \((\varphi\leftrightarrow\psi)\mathrel{:=}((\varphi\to\psi)\land(\psi\to\varphi))\)
  of the defined biconditional in \cref{sec:prop-syntax-semantics-01}.
  Outside those scopes the same macro prints the editorial
  \(\mathrel{:\equiv}\).  In both renderings the left expression abbreviates
  the right one; neither is the object-language identity sign.
  \item \emph{Source alpha-equivalence.}  In
  \cref{sec:fol-syntax-semantics-04}, the source chapters that depend on
  it, and the side condition of the alpha rule in the rule sheet of
  \cref{sec:rule-sheets-01}, \(\varphi\equiv_\alpha\psi\) is
  alpha-equivalence of formulas: the least
  equivalence relation that is a congruence for every constructor and
  contains \(Qx\,\varphi\equiv_\alpha Qy\,(\varphi[y/x])\) whenever
  \(y\notin\operatorname{Var}(\varphi)\).  The subscript is part of the name;
  it is not a variable.
  \item \emph{Schematic register entry.}  The style rules of
  \cref{sec:latex-style-01} and the notation table of \cref{app:notation}
  list \(\varphi\equiv\psi\) as a metalinguistic relation that is available
  only after a semantic class has been declared: it then states that the two
  formulas are logically equivalent with respect to that class.  This entry
  is a schema, not a further relation.  For the class of all valuations it
  is the registered CPL role; for the class of all first-order structures
  with their assignments it is item~1.
  \item \emph{Number-theoretic congruence.}  In the proof of
  \cref{lem:goedel-beta} in \cref{sec:decision-hierarchy-02},
  \(t\equiv a\pmod n\) relates two natural numbers \(t,a\) with respect to a
  modulus \(n\geq1\): it holds iff \(n\) divides the integer difference
  \(t-a\).  The operands are numbers and the modulus is displayed by
  \(\pmod{\cdot}\); the sign is used in the metalanguage of that proof, as
  in the source's appeal to the Chinese remainder theorem for the moduli
  \(1+(i+1)m\), and it is not a formula of the arithmetic object language.
  \item \emph{Many-one equivalence.}  In \cref{rem:jump-convention} and the
  results that cite it, \(A\equiv_m B\) relates two sets of natural numbers:
  it holds iff \(A\leq_m B\) and \(B\leq_m A\).  The subscript \(m\) is part
  of the name, not a modulus and not a variable.
\end{enumerate}
The registered CPL role (\(\varphi\equiv\psi\) for equality of truth values
under every valuation) is the propositional counterpart of item~1.
\end{hbconvention}

These uses meet the editorial defaults in four places.  Item~1 and the
editorial \(\varphi\Leftrightarrow\psi\) name the same relation on
first-order formulas, so a source statement \(\varphi\equiv\psi\) is read
as \(\varphi\Leftrightarrow\psi\) in editorial text, without change of
meaning.  Item~1 and elementary equivalence \(\mathcal M\equiv\mathcal N\)
differ in type: the first relates two formulas over all structures, the
second relates two structures over all sentences; the operands always show
which is meant.  The source also writes elementary equivalence with the same
glyph, as satisfaction of the same \(\mathcal L\)-sentences, in
\cref{thm:upward-ls}.  Item~4 and the registered \(=_\alpha\) are the same
relation, defined by the same closure conditions; the rule sheets and the
editorial text write it with the registered glyph.  The
ordinal-indexed back-and-forth relation that \cref{def:p03-morley-approximants}
writes with the same glyph \(\equiv_\alpha\) is unrelated: there the
subscript is an ordinal and the operands are tuples in structures.  No
statement about one of the two relations transfers to the other.  Item~3 and
the editorial definitional sign differ only in glyph.  Item~6 shares the
bare glyph with items~1, 2 and~5 and with elementary equivalence, but its
operands are numbers and it always carries a displayed modulus, so it cannot
be read as a relation between formulas or structures; the only passage the
source uses from it, to the remainder on division, is part~5 of the next
lemma.  Item~7 and the many-one reducibility preorder \(\leq_m\) are linked by
definition; item~7 relates sets, never formulas.

\begin{hblemma}{thm:p01-w02-p01-025-equiv-bridges}{proved=annex:p01-w02-p01-025-equiv-bridges}{Valid passages between the equivalence notions}
Let \(\varphi,\psi\) be formulas of a first-order language \(\mathcal L\),
let \(\sigma,\tau\) be \(\mathcal L\)-sentences and let \(\mathcal M,\mathcal
N\) be \(\mathcal L\)-structures.
\begin{enumerate}
  \item \(\varphi\equiv\psi\) in the first-order sense iff
  \((\varphi\to\psi)\land(\psi\to\varphi)\) is valid.
  \item If \(\varphi\equiv_\alpha\psi\), then \(\varphi\equiv\psi\) and
  \(c(\varphi)=c(\psi)\).  The converse fails whenever \(\mathcal L\) has a
  unary relation symbol \(P\): \(P(x)\land P(x)\equiv P(x)\), but the two
  formulas are not alpha-equivalent.
  \item If \(\sigma\equiv\tau\) and \(\mathcal M\equiv\mathcal N\), then
  \(\mathcal M\entails\sigma\) iff \(\mathcal N\entails\tau\).
  \item For \(\mathsf{CPL}\) formulas \(\alpha,\beta\), the registered CPL
  relation \(\alpha\equiv\beta\), and hence the metalinguistic sign of
  item~2 of \cref{def:p01-w02-p01-025-equivalence-roles}, holds iff
  \(\alpha\leftrightarrow\beta\) is true under every valuation.
  \item For natural numbers \(t,a,n\) with \(n\geq1\), if
  \(t\equiv a\pmod n\) in the sense of item~6 of
  \cref{def:p01-w02-p01-025-equivalence-roles} and \(a<n\), then \(a\) is
  the remainder of \(t\) on division by \(n\).
\end{enumerate}
\end{hblemma}

\begin{hbexample}{def:p01-w02-p01-025-signs-example}{Reading one line in two notations}
The source normal-form section states \(r(\varphi,V)\equiv\varphi\) for the
rectified variant of \cref{lem:rectified-variant}.  By item~1 of
\cref{def:p01-w02-p01-025-equivalence-roles} this is first-order logical
equivalence, written \(r(\varphi,V)\Leftrightarrow\varphi\) in editorial
notation.  The same lemma also states \(r(\varphi,V)\equiv_\alpha\varphi\),
which is the stronger syntactic claim; part~2 of
\cref{thm:p01-w02-p01-025-equiv-bridges} is the passage from the second
statement to the first.  Reading \(\equiv\) there as elementary equivalence
would be a type error, since \(r(\varphi,V)\) and \(\varphi\) are formulas.
\end{hbexample}
